\documentclass[conference]{IEEEtran}

% --- Paquetes ---
\usepackage[utf8]{inputenc}
\usepackage[T1]{fontenc}
\usepackage[english]{babel}
\usepackage{microtype}
\usepackage{cite}
\usepackage{amsmath,amssymb,amsfonts}
\usepackage{graphicx}
\usepackage{float}        % permite [H] para fijar figuras
\usepackage{textcomp}
\usepackage{xcolor}
\usepackage{listings}
\usepackage{booktabs}     % tablas: \toprule, \midrule, \bottomrule
\usepackage{pgfplots}     % gráficas dibujadas desde los CSV de data/
\pgfplotsset{compat=1.17}
\usepackage{hyperref}

% --- Estilo común de las gráficas ---
% Los CSV de data/ los generan los scripts de ossim/experiments/. Si se
% vuelven a correr, las gráficas se actualizan solas al compilar.
\pgfplotsset{
    expplot/.style={
        width=\columnwidth,
        height=0.62\columnwidth,
        grid=major,
        grid style={gray!25},
        tick label style={font=\footnotesize},
        label style={font=\footnotesize},
        legend style={font=\scriptsize, draw=none, fill=none},
        legend cell align=left,
        mark size=1.6pt,
        line width=0.8pt,
    },
}

\graphicspath{{./}{figures/}}

% --- Hipervínculos ---
\hypersetup{
    colorlinks=true,
    linkcolor=black,
    urlcolor=blue!70!black,
    citecolor=black
}

% --- Estilo de terminal / código ---
\definecolor{codebg}{rgb}{0.96,0.96,0.96}
\definecolor{codegray}{rgb}{0.45,0.45,0.45}
\definecolor{codegreen}{rgb}{0.10,0.45,0.20}
\definecolor{codeblue}{rgb}{0.10,0.25,0.65}

\lstdefinestyle{bashstyle}{
    language=bash,
    backgroundcolor=\color{codebg},
    basicstyle=\ttfamily\footnotesize,
    keywordstyle=\color{codeblue},
    commentstyle=\color{codegreen},
    breaklines=true,
    breakatwhitespace=false,
    columns=fullflexible,
    keepspaces=true,
    numbers=none,
    showspaces=false,
    showstringspaces=false,
    frame=single,
    rulecolor=\color{codegray},
    xleftmargin=2pt,
    xrightmargin=2pt,
    aboveskip=6pt,
    belowskip=6pt
}
\lstset{style=bashstyle}

% --- Marcas de trabajo pendiente ---
% \todo{...} se ve en rojo en el PDF. Antes de entregar, buscar "\todo" y
% "\figplaceholder" en todos los archivos: no debe quedar ninguno.
\newcommand{\todo}[1]{\textcolor{red}{\textbf{[TODO:} #1\textbf{]}}}
% Recuadro que ocupa el lugar de una gráfica que todavía no existe.
\newcommand{\figplaceholder}[1]{%
    \fbox{\parbox[c][0.18\textheight][c]{0.92\columnwidth}{\centering\small #1}}}

% --- Nombre del simulador ---
\newcommand{\ossim}{\texttt{ossim}}

\begin{document}

% --- Título y autores ---
\title{COMP2040: Operating Systems \\ Group Project: \ossim, an Integrated Kernel Simulator}

\author{
    \IEEEauthorblockN{Daniela Méndez Ramírez}
    \IEEEauthorblockA{\textit{College of Engineering and Computer Science} \\
    \textit{VinUniversity}, Hanoi, Vietnam \\
    Student ID: VU2614027}
    \and
    \IEEEauthorblockN{Nguyen Huy An Khanh}
    \IEEEauthorblockA{\textit{College of Engineering and Computer Science} \\
    \textit{VinUniversity}, Hanoi, Vietnam \\
    Student ID: \todo{ID}}
    \and
    \IEEEauthorblockN{\todo{Member 3 name}}
    \IEEEauthorblockA{\textit{College of Engineering and Computer Science} \\
    \textit{VinUniversity}, Hanoi, Vietnam \\
    Student ID: \todo{ID}}
}

\maketitle

% --- Abstract ---
\begin{abstract}
In a real operating system, CPU scheduling, memory management, disk scheduling and synchronization run at the same time and depend on each other. In this project, our team builds \ossim, a kernel simulator written in C in which all of these mechanisms work together. The simulator moves forward one clock tick at a time. It includes six CPU schedulers, deadlock detection, three page-replacement policies, four disk schedulers and a tool that generates random workloads. We check every policy against examples that we solve by hand, and all 45 automatic tests pass. Then we use the simulator for experiments. The first experiment shows that SJF and SRTF give the lowest waiting time, while RR and MLFQ give the lowest response time, at every level of I/O. The second shows that CPU utilization falls from 66\% to 13\% when memory becomes too small, with a clear start of thrashing. The third studies \todo{theme experiment}.
\end{abstract}

\begin{IEEEkeywords}
Operating Systems, CPU Scheduling, Paging, Page Replacement, Disk Scheduling, Deadlock Detection, Simulation.
\end{IEEEkeywords}

% --- Secciones ---
% Un archivo por sección, para que cada miembro pueda editar el suyo sin
% conflictos en git.
% --- I. Introduction ---
\section{Introduction}
An operating system (OS) has many jobs. It decides which program uses the processor, it shares the memory between programs, it sends requests to the disk, and it stops programs from using the same resource in a wrong way. In the labs of this course, we studied each of these jobs alone: a CPU scheduler in one week, a page-replacement algorithm in another week, and a disk scheduler later. In a real system, however, these mechanisms work at the same time and depend on each other~\cite{brief}. For example:
\begin{itemize}
    \item When a program needs a part of its memory that is not loaded (a \emph{page fault}), it must stop using the processor and wait for the disk.
    \item When the disk finishes, the program becomes ready to run again, and the CPU scheduler must decide when it runs.
    \item When a program waits for a lock that another program holds, the scheduler cannot choose it until the lock is free.
\end{itemize}
So a slow disk can make the CPU idle, and a small memory can make the whole system slow. These effects are hard to see when each mechanism is studied alone.

\subsection{Goal of the Project}
The goal of this project is to build one simulator, called \ossim, in which all of these mechanisms run together, and then to use it as a tool for experiments. A \emph{simulator} is a program that imitates the behavior of another system. \ossim{} does not run real programs. Instead, it reads a text file (a \emph{workload}) that describes a set of imaginary processes: when they arrive, how long they use the CPU, when they do input/output (I/O), which pages of memory they use, and which locks they take. Then it simulates what an OS would do with them and measures the result.

The command line of \ossim{} selects the policies to use: the CPU scheduler, the page-replacement policy, the disk scheduler, the number of memory frames and the disk parameters. In this way, the same workload can be run with different policies and the results can be compared fairly.

\subsection{What the Team Built}
The course gave a starting version of \ossim{} (the \emph{scaffold}). It already had the main loop of the simulator, the simplest policies (FCFS and Round Robin for the CPU, FIFO for memory and FCFS for the disk) and the code for semaphores and resources. Our team completed the rest in two checkpoints:
\begin{itemize}
    \item \textbf{Checkpoint 1} (processes, scheduling and synchronization): four more CPU schedulers (SJF, SRTF, Priority and MLFQ) and deadlock detection.
    \item \textbf{Checkpoint 2} (memory and storage): two more page-replacement policies (LRU and Clock), three more disk schedulers (SSTF, SCAN and C-LOOK) and a workload generator.
\end{itemize}
For every policy, we first solved an example by hand and then checked that the simulator gives exactly the same answer. After that, we used the simulator for experiments, and we extended it with one extra mechanism (the \emph{theme}).

\subsection{Structure of this Report}
Section~\ref{sec:background} explains the OS concepts that the rest of the report uses. Section~\ref{sec:design} describes how the simulator is built. Section~\ref{sec:validation} shows how we checked it against hand-computed examples. Section~\ref{sec:experiments} presents two experiments, and Section~\ref{sec:theme} presents the theme and its experiment. Section~\ref{sec:conclusion} summarizes the results. All the code is available in the team repository~\cite{repo}.

% --- II. Background ---
% Conceptos que el resto del reporte usa. Cada término nuevo va en negritas
% la primera vez que aparece, con una explicación corta.
\section{Background}
\label{sec:background}
This section explains the main concepts of the project in simple terms. All of them come from the course textbook~\cite{silberschatz}.

\subsection{Processes and Their States}
A \textbf{process} is a program in execution. The OS keeps the information of each process in a \textbf{process control block (PCB)}: its identifier (\emph{pid}), its state, its priority, its memory information and some counters.

A process is always in one of five \textbf{states}:
\begin{itemize}
    \item \textbf{New}: the process is being created.
    \item \textbf{Ready}: the process can run and waits for the CPU in the \textbf{ready queue}.
    \item \textbf{Running}: the process is using the CPU. With one CPU, only one process can be running.
    \item \textbf{Waiting} (or blocked): the process waits for something, for example the end of an I/O operation, a page from the disk or a lock.
    \item \textbf{Terminated}: the process has finished.
\end{itemize}

The life of a process is a sequence of \textbf{CPU bursts} (periods of computing) and \textbf{I/O bursts} (periods of waiting for a device). A process with long CPU bursts is \emph{CPU-bound}; a process with many short bursts and a lot of I/O is \emph{I/O-bound}.

\subsection{CPU Scheduling}
The \textbf{CPU scheduler} chooses which ready process runs next. When the CPU changes from one process to another, the OS makes a \textbf{context switch}: it saves the state of the old process and loads the state of the new one. A context switch is pure overhead, because no useful work is done during it.

A scheduler is \textbf{preemptive} if it can take the CPU away from a running process before its burst ends, for example when a more important process arrives. A \textbf{non-preemptive} scheduler lets the process run until it finishes its burst or blocks. Some preemptive schedulers give each process a \textbf{time quantum} $q$, a maximum number of ticks that it can run before another process gets a turn.

The six schedulers of \ossim{} are:
\begin{itemize}
    \item \textbf{FCFS} (First-Come, First-Served): processes run in the order they become ready.
    \item \textbf{RR} (Round Robin): like FCFS, but each process runs for at most one quantum and then goes to the end of the queue.
    \item \textbf{SJF} (Shortest Job First): the process with the shortest next CPU burst runs first. It is non-preemptive.
    \item \textbf{SRTF} (Shortest Remaining Time First): the preemptive version of SJF. If a new process has less work left than the running one, it takes the CPU.
    \item \textbf{PRIO} (Priority): the process with the highest priority runs first. In \ossim, a smaller number means a higher priority.
    \item \textbf{MLFQ} (Multilevel Feedback Queue): several queues with different priorities. A process that uses its whole quantum moves to a lower queue, so long CPU-bound processes go down and short interactive processes stay at the top.
\end{itemize}

To compare schedulers, we use these \textbf{metrics} for each process. Lower values are better for all of them.
\begin{itemize}
    \item \textbf{Turnaround time}: time from arrival to termination.
    \item \textbf{Waiting time}: total time spent in the ready queue.
    \item \textbf{Response time}: time from arrival until the process runs for the first time. It matters for interactive programs.
\end{itemize}
For the whole system, we also use \textbf{CPU utilization}, the percentage of time the CPU does useful work, and \textbf{throughput}, the number of processes that finish per unit of time.

\subsection{Synchronization and Deadlock}
When processes share data or devices, they must coordinate. A \textbf{semaphore} is an integer counter with a list of waiting processes and two operations. \texttt{wait} decreases the counter if it is greater than zero; if it is zero, the process blocks and joins the list. \texttt{signal} wakes up the first process in the list or, if nobody is waiting, increases the counter. A semaphore with initial value 1 works as a \textbf{lock}: only one process at a time can be inside the protected part of the code.

A \textbf{resource} is something that processes request and release, such as a printer. A resource can have several identical \textbf{instances}. A \textbf{deadlock} happens when a group of processes wait for each other forever. For example, process A holds resource R1 and waits for R2, while process B holds R2 and waits for R1. Neither of them can continue. \textbf{Deadlock detection} is an algorithm that looks at who holds what and who waits for what, and finds the processes that can never continue.

\subsection{Memory and Paging}
Each process sees its memory as a list of fixed-size blocks called \textbf{pages}. The physical memory is divided into blocks of the same size called \textbf{frames}. The \textbf{page table} of a process says, for each page, in which frame it is, or that it is not in memory. This technique is called \textbf{paging}.

When a process uses (\emph{references}) a page that is not in memory, a \textbf{page fault} happens. The process blocks, the OS reads the page from the disk into a free frame, and then the process can continue. If there is no free frame, the OS must choose a page to remove (the \textbf{victim}). This choice is the job of the \textbf{page-replacement policy}:
\begin{itemize}
    \item \textbf{FIFO}: remove the page that was loaded first.
    \item \textbf{LRU} (Least Recently Used): remove the page that has not been used for the longest time.
    \item \textbf{Clock} (or second chance): a cheaper approximation of LRU. Each frame has a \textbf{reference bit} that is set when its page is used. A pointer (the \emph{hand}) moves around the frames. A page with the bit set gets a second chance: the bit is cleared and the hand moves on. The first page with the bit clear is the victim.
\end{itemize}
In \ossim, replacement is \textbf{global}: all processes share the same frames, so one process can take a frame from another.

Programs usually show \textbf{locality of reference}: in a short period they use only a small set of pages, called their \textbf{working set}. If the frames cannot hold the working sets of all processes, the processes keep removing each other's pages. Then almost every reference is a page fault, the processes spend most of the time waiting for the disk, and CPU utilization falls. This situation is called \textbf{thrashing}.

\subsection{Disk Scheduling}
A hard disk stores data on circular tracks; the tracks at the same position on all surfaces form a \textbf{cylinder}. A \textbf{head} must move to the correct cylinder before it can read or write. The time of this movement is the \textbf{seek time}, and it grows with the distance. The time to move the data itself is the \textbf{transfer time}. Since seeking is slow, the order in which requests are served matters. We measure it with the \textbf{total head movement}, the number of cylinders the head crosses. The four disk schedulers of \ossim{} are:
\begin{itemize}
    \item \textbf{FCFS}: serve the requests in arrival order.
    \item \textbf{SSTF} (Shortest Seek Time First): serve the request closest to the head.
    \item \textbf{SCAN} (the elevator algorithm): move in one direction serving requests until the end of the disk, then turn around.
    \item \textbf{C-LOOK}: move in one direction only as far as the last request, then jump back to the first request at the other end and continue in the same direction.
\end{itemize}

\subsection{Tick-Driven Simulation}
\ossim{} is a \textbf{tick-driven} simulator. Time is divided into small, equal steps called \textbf{ticks}, and the simulator updates the whole system once per tick. An event that will happen in the future, like the end of an I/O operation, is stored with its time and handled when the clock reaches it. Because the simulator has no real randomness, the same input always gives the same output. We call this property \textbf{determinism}, and it is what makes hand-computed tests possible.

% --- III. Simulator Design ---
\section{Simulator Design}
\label{sec:design}
This section describes how \ossim{} is built: how the code is organized, what the input and the output look like, how one tick works, and how each of the seven algorithms is implemented.

\subsection{Code Organization}
\ossim{} is written in standard C (C99) and uses no external libraries. It is built with \texttt{make}, using the flags \texttt{-std=c99 -Wall -Wextra -Werror}. With \texttt{-Werror}, every compiler warning is treated as an error, so the code must compile without any warning. Each part of the simulator is a \textbf{module}: a header file in \texttt{include/} that declares what the module offers, and a source file in \texttt{src/} that contains the code. Table~\ref{tab:modules} lists the modules used in Checkpoint~1.

\begin{table}[H]
\centering
\caption{Modules of \ossim{} used in Checkpoint 1.}
\label{tab:modules}
\begin{tabular}{p{0.27\columnwidth}p{0.63\columnwidth}}
\toprule
File & Purpose \\
\midrule
\texttt{main.c} & Reads the command line and starts the simulation. \\
\texttt{workload.c} & Reads the workload file and creates the processes. \\
\texttt{kernel.c} & The tick loop and the life of every process. \\
\texttt{event.c} & A list of future events, sorted by time. \\
\texttt{queue.c} & Queues of processes (ready queues, waiting lists). \\
\texttt{sched.c} & The six CPU schedulers. \\
\texttt{sync.c} & Semaphores, resources and deadlock detection. \\
\texttt{metrics.c} & The final report and the Gantt chart. \\
\texttt{tests/run.sh} & The automatic tests. \\
\bottomrule
\end{tabular}
\end{table}

The simulator also contains modules for memory and the disk, which belong to the next checkpoint. In this report, all runs use the option \texttt{--no-paging}, which turns off page faults, so the results depend only on the CPU scheduler and on synchronization.

\subsection{Input: the Workload File}
A workload file describes semaphores, resources and processes, one instruction per line. Each process has an arrival time, a priority and a list of operations that it does in order. The example below shows a process that computes, waits for I/O, and then uses a lock:

\begin{lstlisting}
semaphore lock 1

process editor arrive=0 prio=2
    cpu 5
    io 3
    wait lock
    cpu 2
    signal lock
end
\end{lstlisting}

The operations are:
\begin{itemize}
    \item \texttt{cpu N}: use the CPU for $N$ ticks.
    \item \texttt{io N}: wait $N$ ticks for a device.
    \item \texttt{wait S}, \texttt{signal S}: semaphore operations.
    \item \texttt{acquire R n}, \texttt{release R n}: take or give back $n$ instances of resource $R$.
\end{itemize}

\subsection{Output: the Report}
At the end of a run, \ossim{} prints one line per process and a summary. The output for the textbook Round Robin example (three processes with bursts 24, 3 and 3, quantum 4) is shown below in a shortened form:

\begin{lstlisting}
$ ./ossim workloads/basic.txt --sched rr --quantum 4 --no-paging --gantt
CPU timeline:
  [0-3] P1 | [4-6] P2 | [7-9] P3 | [10-29] P1 |
summary
  ticks=30 finished=3/3 aborted=0
  throughput=0.1000 procs/tick context_switches=4
  avg_turnaround=15.67 avg_waiting=5.67 avg_response=3.67
\end{lstlisting}

The option \texttt{--gantt} prints the Gantt chart. Each interval \texttt{[a-b]} includes both ticks, so \texttt{[0-3] P1} means that P1 ran in ticks 0, 1, 2 and 3. The option \texttt{--trace} prints every change of state, which is very useful to follow a hand-computed example step by step.

\subsection{The Tick Loop}
The center of the simulator is the \textbf{tick loop} in \texttt{kernel.c}. In every tick, the kernel does five steps:
\begin{enumerate}
    \item \textbf{Deliver events.} Events whose time has come are handled: new processes arrive and I/O operations end. A process whose wait has ended goes back to the ready queue.
    \item \textbf{Schedule.} The kernel asks the scheduler if the running process must be preempted. If the CPU is free, it asks the scheduler for the next process and gives it the CPU (\emph{dispatch}).
    \item \textbf{Account.} Every ready process adds one tick to its waiting time, and every blocked process adds one tick to its blocked time.
    \item \textbf{Execute.} The running process uses one tick of its current CPU burst. When the burst ends, the process moves on to its next operation.
    \item \textbf{Check and advance.} If deadlock detection is enabled and it is time, it runs. Then the clock moves forward by one tick.
\end{enumerate}
Operations that take no time, like \texttt{signal} or a \texttt{wait} that does not block, are done immediately when a process reaches them. As a result, a process is in the ready queue only when its next operation is a CPU burst. The kernel is the only part of the code that changes the state of a process; the schedulers only manage their ready queues.

\subsection{Schedulers as Interchangeable Parts}
The kernel does not know which scheduler is running. Each scheduler is a C structure that contains \textbf{function pointers}, which are variables that store the address of a function. The kernel calls the functions through these pointers, so changing the scheduler only means choosing a different structure. Every scheduler provides the same functions:
\begin{itemize}
    \item \texttt{enqueue(p)}: process $p$ has become ready.
    \item \texttt{should\_preempt(r)}: should the running process $r$ leave the CPU now?
    \item \texttt{pick()}: remove and return the next process to run.
    \item \texttt{on\_dispatch(p)}, \texttt{on\_tick(r)}: update counters, such as the part of the quantum already used.
\end{itemize}
Thanks to this design, we could add four schedulers without changing the kernel. The ready queues use two helper functions: \texttt{q\_pop\_min()} removes the process with the smallest value of a \emph{key} (for example, the remaining burst), and \texttt{q\_min\_key()} returns that smallest value without removing anything. When two processes have the same key, the one that entered the queue first wins. This rule makes every run deterministic, which is necessary for the hand-computed tests.

\subsection{The Six CPU Schedulers}
\textbf{FCFS} (given) keeps one first-in, first-out queue and never preempts.

\textbf{RR} (given) uses the same queue, but counts the ticks that the running process has used. When the quantum is used up and another process is ready, the running process goes to the end of the queue. If no other process is ready, it keeps the CPU and starts a new quantum.

\textbf{SJF} keeps one ready queue and, when the CPU is free, takes the process with the shortest remaining time in its \emph{current} burst. It never preempts. In a real OS, the length of the next burst is unknown and must be predicted, for example with exponential averaging~\cite[Sec.~5.3.2]{silberschatz}. The simulator knows it exactly, which makes SJF an ideal reference.

\textbf{SRTF} uses the same key. Before every tick, it compares the running process with the best ready process. It preempts only when the ready process has \emph{strictly} less time left, so two equal values do not cause a useless context switch.

\textbf{PRIO} takes the ready process with the smallest priority number. A process with a higher priority preempts the running one at once. Processes with the same priority share the CPU in Round Robin with quantum $q$. PRIO has no aging, so a process with a low priority can starve.

\textbf{MLFQ} has three queues. Level 0 has quantum $q$, level 1 has $2q$ and level 2 has $4q$. A new process starts at level 0. The scheduler always takes a process from the highest non-empty level, and a process at a higher level preempts a process at a lower level. A process that uses its whole quantum moves down one level; a process that blocks before the end of its quantum keeps its level. Every 20 ticks, a \textbf{priority boost} moves all processes back to level 0. Without the boost, a long process at level 2 could starve behind new short processes.

\subsection{Deadlock Detection}
Each resource keeps three pieces of information: how many instances are free (\textit{Available}), how many each process holds (\textit{Allocation}) and how many each blocked process is asking for (\textit{Request}). Deadlock detection uses the algorithm for resources with several instances~\cite[Sec.~8.6.2]{silberschatz}:
\begin{enumerate}
    \item Copy \textit{Available} into a work vector \textit{Work}. No process is marked.
    \item Find an unmarked process whose \textit{Request} is not larger than \textit{Work}. This process could get what it asks for, finish, and give back everything it holds, so add its \textit{Allocation} to \textit{Work} and mark it. Repeat until no more processes can be marked.
    \item The blocked processes that are still unmarked are deadlocked.
\end{enumerate}
The idea of step 2 is optimistic: if a process can finish, it will release its resources and help the others. A process that cannot be marked has no possible way to get what it needs.

The option \texttt{--deadlock-every N} runs the check every $N$ ticks. The check only reads the state of the system and never changes it. When it finds a deadlock, it prints a line such as \texttt{DEADLOCK t=15 pids: 1 2}, with the process identifiers from small to large. Since no process is terminated to break the deadlock, the same group would be found at every check, so it is printed only once.

\subsection{Testing}
The file \texttt{tests/run.sh} contains the automatic tests, and \texttt{make test} runs them. Each test runs \ossim{} on a workload and checks one result against a value that we computed by hand. Some tests check a number of the summary, for example \texttt{avg\_waiting=7.00}; others check a piece of the Gantt chart, for example \texttt{[16-16] P5 | [17-19] P3}. A test prints PASS if the result is correct, FAIL if it is wrong, and SKIP if the algorithm is not implemented yet. In this way, the tests were written before the code, and they showed our progress during the checkpoint.

\subsection{Teamwork}
The team works in a shared git repository~\cite{repo}. Each member makes regular commits, and the design note \texttt{CHECKPOINT1.md} records the rules of the algorithms, the hand-computed examples and the results.

% --- IV. Validation ---
\section{Validation}
\label{sec:validation}
A simulator is only useful if its results are correct. To check the seven algorithms, we use \textbf{hand-computed cases}: small examples that we solve on paper, step by step, before we run the simulator. If the simulator gives exactly the same numbers and the same Gantt chart, we can trust the algorithm. Some examples come from the textbook~\cite{silberschatz}, so their answers are already known; when the textbook has no answer, we made our own example. All these checks are in \texttt{tests/run.sh}. The file has 36 checks for the given scaffold and for this checkpoint, and all of them pass.

\subsection{A Worked Example: SRTF}
To show how a case is computed by hand, we use the SRTF example of the textbook~\cite[Sec.~5.3.2]{silberschatz}. Four processes arrive at times 0, 1, 2 and 3 with bursts of 8, 4, 9 and 5 ticks:
\begin{itemize}
    \item At $t=0$, only P1 is ready, so it runs.
    \item At $t=1$, P2 arrives with 4 ticks of work, while P1 has 7 left. Since $4 < 7$, P2 preempts P1.
    \item P2 runs until $t=5$. P3 (9 ticks) and P4 (5 ticks) arrive meanwhile, but neither has less work left than P2.
    \item At $t=5$, the ready processes have 7 (P1), 9 (P3) and 5 (P4) ticks left. P4 runs until $t=10$, then P1 until $t=17$, and finally P3 until $t=26$.
\end{itemize}
The waiting times are 9, 0, 15 and 2 ticks, so the average is $26/4 = 6.50$. There are 5 context switches (P1, P2, P4, P1, P3). The simulator prints \texttt{avg\_waiting=6.50} and \texttt{context\_switches=5}, and its Gantt chart is identical.

\subsection{Textbook Cases}
Table~\ref{tab:cpu_textbook} shows all the textbook examples that we use~\cite[Sec.~5.3]{silberschatz}. In every case, the simulator gives exactly the textbook value.

\begin{table}[H]
\centering
\caption{Average waiting time on the textbook examples.}
\label{tab:cpu_textbook}
\begin{tabular}{llr}
\toprule
Policy & Workload & Avg. waiting \\
\midrule
FCFS & bursts 24, 3, 3 & 17.00 \\
RR ($q=4$) & bursts 24, 3, 3 & 5.67 \\
SJF & bursts 6, 8, 7, 3 & 7.00 \\
SRTF & arrivals 0--3, bursts 8, 4, 9, 5 & 6.50 \\
\bottomrule
\end{tabular}
\end{table}

\subsection{All Six Schedulers on One Workload}
The textbook examples are simple: the processes do no I/O and have no priorities. To test the schedulers in a harder situation, we use the workload \texttt{arrivals.txt}. It has six jobs (J1 to J6) with different arrival times and priorities, and two of them do I/O between two CPU bursts:

\begin{table}[H]
\centering
\caption{The workload \texttt{arrivals.txt} (priority 1 is the highest).}
\label{tab:arrivals_input}
\begin{tabular}{lrrl}
\toprule
Job & Arrival & Priority & Operations \\
\midrule
J1 & 0 & 2 & cpu 7, io 3, cpu 4 \\
J2 & 2 & 1 & cpu 9 \\
J3 & 4 & 3 & cpu 3, io 5, cpu 3 \\
J4 & 6 & 2 & cpu 2 \\
J5 & 8 & 3 & cpu 5 \\
J6 & 12 & 1 & cpu 6 \\
\bottomrule
\end{tabular}
\end{table}

We drew the Gantt chart of every scheduler by hand (Appendix~\ref{app:gantt}) and compared it with the output of \texttt{--gantt}. They match tick by tick. Table~\ref{tab:arrivals} shows the averages.

\begin{table}[H]
\centering
\caption{Results on \texttt{arrivals.txt} (hand-computed and simulated).}
\label{tab:arrivals}
\begin{tabular}{lrrrr}
\toprule
Policy & Waiting & Response & Turnaround & Switches \\
\midrule
FCFS & 14.83 & 10.17 & 22.67 & 8 \\
RR ($q=2$) & 18.33 & 3.17 & 26.17 & 22 \\
SJF & 10.00 & 9.00 & 17.83 & 8 \\
SRTF & 9.83 & 9.00 & 17.67 & 9 \\
PRIO ($q=2$) & 12.83 & 8.33 & 20.67 & 12 \\
MLFQ ($q=2$) & 15.00 & 0.00 & 22.83 & 18 \\
\bottomrule
\end{tabular}
\end{table}

The table already shows the typical behavior of each algorithm:
\begin{itemize}
    \item \textbf{SRTF} has the lowest waiting time, as the theory says. The only difference with SJF happens at $t=17$: J3 comes back from I/O with 3 ticks of work, while J5 still has 4, so SRTF gives the CPU to J3. At $t=4$, J3 arrives with 3 ticks and J1 also has 3 left; since the values are equal, J1 keeps the CPU. Both moments are checked by the tests.
    \item \textbf{RR} has the best response time after MLFQ, but the worst waiting time and 22 context switches, because the quantum is short and every job is cut many times.
    \item \textbf{PRIO} lets J2 preempt J1 at $t=2$, because J2 has priority 1. The jobs with priority 3 (J3 and J5) wait the longest.
    \item \textbf{MLFQ} gives the CPU to every job as soon as it arrives (response time 0), because new jobs enter at level 0, above the jobs that have already used their quantum. The boost at $t=20$ is also checked: it moves J5 back to level 0 and gives it a new quantum.
\end{itemize}

\subsection{Equal Priorities}
The workload \texttt{prio\_ties.txt} has two processes with the same priority and bursts of 5 ticks. With $q=2$, PRIO must alternate between them like Round Robin: P1 runs ticks 0--1, P2 ticks 2--3, P1 ticks 4--5, and so on. The simulator gives this exact order, an average waiting time of 4.50 and an average response time of 1.00, as computed by hand.

\subsection{Deadlock Detection}
We check the algorithm on two workloads. The first, \texttt{deadlock.txt}, has the classic case with one instance per resource: process A holds R1 and wants R2, and process B holds R2 and wants R1. A third process, C, does not use these resources. With \texttt{--deadlock-every 5}, the simulator prints \texttt{DEADLOCK t=15 pids: 1 2} at the first check after both processes block, and C finishes normally.

The second workload, \texttt{deadlock\_multi.txt}, tests the multi-instance case. Resources R1 and R2 have two instances each. Process A takes both instances of R1 and then asks for one instance of R2; process B takes both instances of R2 and then asks for one instance of R1. By hand, at $t=1$:
\begin{itemize}
    \item \textit{Available} $= (0, 0)$, because all instances are held.
    \item \textit{Allocation}: A $= (2, 0)$, B $= (0, 2)$.
    \item \textit{Request}: A $= (0, 1)$, B $= (1, 0)$.
    \item \textit{Work} $= (0, 0)$. The request of A does not fit, and the request of B does not fit either, so no process can be marked.
\end{itemize}
Both processes stay unmarked, so both are deadlocked. The simulator prints \texttt{DEADLOCK t=1 pids: 1 2}, which matches.

Without deadlock detection, \texttt{deadlock.txt} must still stop with A and B unfinished and C finished, and the producer-consumer workload, which uses three semaphores, must finish correctly. We run the producer-consumer workload with all six schedulers, and it finishes in every case.

% --- V. Experiments ---
% El enunciado pide al menos tres experimentos, cada uno con: pregunta,
% método, gráfica o tabla, y discusión. El experimento del tema cuenta como
% uno de los tres (está en 06-theme.tex).
% Las gráficas leen data/*.csv, generados por ossim/experiments/*.sh.
\section{Experiments}
\label{sec:experiments}
In this section, we use the simulator as an instrument to answer questions about the behavior of the system. Each experiment has a question, a method, a plot and a discussion. All experiments use workloads made by \texttt{wlgen} with seeds 1 to 10, so they can be repeated exactly. With a single workload, a result could depend on chance, so every value is the average over the ten seeds. The scripts are in \texttt{ossim/experiments/}, and the commands are in Appendix~\ref{app:commands}.

% ------------------------------------------------------------------
\subsection{Experiment 1: Schedulers under Varying I/O Intensity}

\textbf{Question.} How do the waiting time and the response time of the six schedulers change when processes do more I/O?

\textbf{Method.} Each workload has 10 processes with 4 CPU bursts of 6 ticks on average, arriving every 3 ticks on average. The I/O intensity $P$ is the percentage of gaps between two bursts in which the process blocks; half of these blocks are disk requests and half are generic I/O of 6 ticks on average. $P$ takes the values 0, 25, 50, 75 and 100. Every scheduler runs every workload with $q = 4$ and \texttt{--no-paging}, so only the CPU scheduler changes between runs.

\textbf{Results.} Fig.~\ref{fig:exp1_wait} shows the average waiting time and Fig.~\ref{fig:exp1_resp} the average response time.

\begin{figure}[H]
    \centering
    \begin{tikzpicture}
    \begin{axis}[expplot,
        xlabel={I/O intensity $P$ (\%)},
        ylabel={Avg. waiting time (ticks)},
        xtick={0,25,50,75,100},
        legend columns=3,
        legend style={at={(0.5,1.03)}, anchor=south},
    ]
        \addplot[blue,   mark=*]        table[x=io, y=fcfs, col sep=comma] {data/exp1_waiting.csv};
        \addplot[red,    mark=square*]  table[x=io, y=rr,   col sep=comma] {data/exp1_waiting.csv};
        \addplot[green!50!black, mark=triangle*] table[x=io, y=sjf,  col sep=comma] {data/exp1_waiting.csv};
        \addplot[orange, mark=diamond*] table[x=io, y=srtf, col sep=comma] {data/exp1_waiting.csv};
        \addplot[violet, mark=pentagon*] table[x=io, y=prio, col sep=comma] {data/exp1_waiting.csv};
        \addplot[gray,   mark=o]        table[x=io, y=mlfq, col sep=comma] {data/exp1_waiting.csv};
        \legend{FCFS, RR, SJF, SRTF, PRIO, MLFQ}
    \end{axis}
    \end{tikzpicture}
    \caption{Average waiting time of the six schedulers as I/O intensity grows.}
    \label{fig:exp1_wait}
\end{figure}

\begin{figure}[H]
    \centering
    \begin{tikzpicture}
    \begin{axis}[expplot,
        xlabel={I/O intensity $P$ (\%)},
        ylabel={Avg. response time (ticks)},
        xtick={0,25,50,75,100},
        legend columns=3,
        legend style={at={(0.5,1.03)}, anchor=south},
    ]
        \addplot[blue,   mark=*]        table[x=io, y=fcfs, col sep=comma] {data/exp1_response.csv};
        \addplot[red,    mark=square*]  table[x=io, y=rr,   col sep=comma] {data/exp1_response.csv};
        \addplot[green!50!black, mark=triangle*] table[x=io, y=sjf,  col sep=comma] {data/exp1_response.csv};
        \addplot[orange, mark=diamond*] table[x=io, y=srtf, col sep=comma] {data/exp1_response.csv};
        \addplot[violet, mark=pentagon*] table[x=io, y=prio, col sep=comma] {data/exp1_response.csv};
        \addplot[gray,   mark=o]        table[x=io, y=mlfq, col sep=comma] {data/exp1_response.csv};
        \legend{FCFS, RR, SJF, SRTF, PRIO, MLFQ}
    \end{axis}
    \end{tikzpicture}
    \caption{Average response time of the six schedulers as I/O intensity grows.}
    \label{fig:exp1_resp}
\end{figure}

\textbf{Discussion.} The results show a clear trade-off between waiting time and response time:
\begin{itemize}
    \item \textbf{Waiting time.} SJF and SRTF have the lowest waiting time at every level of I/O (about 86 to 110 ticks), as the theory predicts~\cite[Sec.~5.3.2]{silberschatz}. RR and MLFQ have the highest (about 145 to 160 ticks), because time slicing keeps every process in the ready queue for longer before it finishes.
    \item \textbf{FCFS and I/O.} The waiting time of FCFS grows from 93 to 148 ticks as $P$ grows. Every time a process comes back from I/O, it joins the end of the queue and may wait behind a long burst again. This is the convoy effect.
    \item \textbf{Response time.} RR and MLFQ have the lowest response time (about 9 to 16 ticks), because a new process gets the CPU after at most a few quanta. With little I/O, FCFS and SJF have the highest response time (93 and 86 ticks at $P=0$). With a lot of I/O, PRIO has the highest (42 ticks at $P=100$), because a process with a low priority can wait a long time for its first turn.
    \item \textbf{Effect of I/O on FCFS response.} The response time of FCFS falls from 93 to 21 ticks as $P$ grows. With more I/O, the running process blocks sooner and frees the CPU, so new processes start earlier.
\end{itemize}
In summary, no scheduler is best for both metrics. SJF and SRTF are best for throughput-oriented work, and RR and MLFQ are best for interactive work. MLFQ gets a response time as low as RR (even slightly lower) without any knowledge of burst lengths.

% ------------------------------------------------------------------
\subsection{Experiment 2: Page Faults and the Onset of Thrashing}

\textbf{Question.} How do the page-fault rate and the CPU utilization change as the number of frames shrinks, and when does thrashing start?

\textbf{Method.} Each workload has 8 processes with 16 pages each, so the total address space is 128 pages. Every process has 4 CPU bursts of 30 ticks on average, and only 10\% of the gaps between bursts block on I/O, so most blocking comes from page faults. The references follow a random walk: 70\% of them move at most one page from the previous one, and the rest jump to any page. The number of frames goes from 128 down to 4, with FIFO, LRU and Clock. All runs use RR ($q=4$) and the FCFS disk, and stop after 20\,000 ticks if they have not finished.

The fault rate is the number of page faults per executed reference. It can be larger than 1: when a newly loaded page is evicted before its owner uses it, the same reference faults again.

\textbf{Results.} Fig.~\ref{fig:exp2_util} shows the CPU utilization (ticks of user work over all ticks) and Fig.~\ref{fig:exp2_fault} the fault rate. The x axis is logarithmic.

\begin{figure}[H]
    \centering
    \begin{tikzpicture}
    \begin{axis}[expplot,
        xmode=log, log basis x=2,
        xlabel={Frames},
        ylabel={CPU utilization (\%)},
        xtick={4,8,16,32,64,128},
        xticklabels={4,8,16,32,64,128},
        ymin=0,
        legend pos=north west,
    ]
        \addplot[blue,   mark=*]        table[x=frames, y=fifo,  col sep=comma] {data/exp2_util.csv};
        \addplot[red,    mark=square*]  table[x=frames, y=lru,   col sep=comma] {data/exp2_util.csv};
        \addplot[green!50!black, mark=triangle*] table[x=frames, y=clock, col sep=comma] {data/exp2_util.csv};
        \legend{FIFO, LRU, Clock}
    \end{axis}
    \end{tikzpicture}
    \caption{CPU utilization as the number of frames shrinks.}
    \label{fig:exp2_util}
\end{figure}

\begin{figure}[H]
    \centering
    \begin{tikzpicture}
    \begin{axis}[expplot,
        xmode=log, log basis x=2,
        ymode=log,
        xlabel={Frames},
        ylabel={Page faults per reference},
        xtick={4,8,16,32,64,128},
        xticklabels={4,8,16,32,64,128},
        legend pos=north east,
    ]
        \addplot[blue,   mark=*]        table[x=frames, y=fifo,  col sep=comma] {data/exp2_faultrate.csv};
        \addplot[red,    mark=square*]  table[x=frames, y=lru,   col sep=comma] {data/exp2_faultrate.csv};
        \addplot[green!50!black, mark=triangle*] table[x=frames, y=clock, col sep=comma] {data/exp2_faultrate.csv};
        \legend{FIFO, LRU, Clock}
    \end{axis}
    \end{tikzpicture}
    \caption{Page-fault rate as the number of frames shrinks (log scale on both axes).}
    \label{fig:exp2_fault}
\end{figure}

\textbf{Discussion.}
\begin{itemize}
    \item \textbf{Enough memory.} With 128 frames every page fits, and the only faults are the first reference to each page (compulsory faults). The fault rate is 0.13 and the utilization is 66\%.
    \item \textbf{Onset of thrashing.} The largest drop happens between 96 and 64 frames: the utilization falls from about 62\% to about 44\%, and the fault rate goes from 0.17 to 0.28. From 32 frames down, the processes spend most of their time waiting for the disk, and the utilization is about 28\% or less.
    \item \textbf{Collapse.} With 4 frames the fault rate jumps to 2.6, so most references fault more than once, and 9\% of the processes do not finish within 20\,000 ticks. Here a page is often evicted by another process before its owner can use it. This is thrashing in its extreme form~\cite[Sec.~10.6]{silberschatz}.
    \item \textbf{Policies.} FIFO, LRU and Clock give almost the same results. The random-walk references of the generator have little reuse that LRU can exploit better than FIFO. On the textbook string, which has more reuse, LRU is clearly better (12 faults against 15). When memory is too small, no replacement policy can help: the problem is the number of processes, not the choice of victim. This is the motivation for load control with the working-set model.
\end{itemize}

% --- VI. Theme Extension ---
% Requisitos del enunciado para el tema:
%   - se activa con una opción de línea de comandos; sin ella, todo igual
%   - al menos una prueba con resultado calculado a mano
%   - un experimento (cuenta como uno de los tres del reporte)
\section{Theme: \todo{theme name}}
\label{sec:theme}
After the core simulator was complete, each team extended it with one extra mechanism, called the \emph{theme}, and used an experiment to study what this mechanism really changes~\cite{brief}. The extension is turned on with a command-line option, so a run without the option behaves exactly as before. This keeps all the results of Sections~\ref{sec:validation} and~\ref{sec:experiments} valid.

\subsection{Motivation}
\todo{What problem does the mechanism solve? One paragraph, with a reference to the textbook.}

\subsection{Design}
\todo{What was added to the simulator, in which files, and the command-line option that turns it on. A run without the option behaves exactly as before.}

\begin{lstlisting}
$ ./ossim workload.txt ... --TODO-option
\end{lstlisting}

\subsection{Validation}
\todo{The hand-computed test: the input, the expected result worked out by hand, and the result of the simulator.}

\subsection{Experiment 3: \todo{title}}

\textbf{Question.} \todo{The question the theme experiment answers.}

\textbf{Method.} \todo{Workloads, parameters, what is varied.}

\textbf{Results.} \todo{Describe the plot or table.}

\begin{figure}[H]
    \centering
    \figplaceholder{Plot or table of the theme experiment.}
    \caption{\todo{caption}}
    \label{fig:exp3}
\end{figure}

\textbf{Discussion.} \todo{What the mechanism actually changes, and when it does not help.}

% --- VII. Conclusion ---
\section{Conclusion}
\label{sec:conclusion}
In this project, our team built \ossim, a kernel simulator in which CPU scheduling, synchronization, paging and disk scheduling run together. Starting from the scaffold, we added four CPU schedulers, deadlock detection, two page-replacement policies, three disk schedulers and a workload generator. We checked every policy against examples solved by hand, and all 45 automatic tests pass.

The experiments show two classic trade-offs. For CPU scheduling, SJF and SRTF give the lowest waiting time, while RR and MLFQ give the lowest response time, at every level of I/O. For memory, CPU utilization drops sharply once the frames can no longer hold the pages that the processes are using, and the choice of replacement policy matters much less than the number of processes in memory. \todo{One sentence with the main finding of the theme experiment.}

The simulator also has some limitations, which are possible future work:
\begin{itemize}
    \item SJF and SRTF know the exact length of the next burst. A real OS must predict it, for example with exponential averaging~\cite[Sec.~5.3.2]{silberschatz}.
    \item PRIO has no \emph{aging}, so a process with a low priority can starve.
    \item A page that has just been loaded is not protected until its owner uses it. With very little memory, this can make the system almost stop, as Experiment~2 shows.
\end{itemize}

% --- Team Contributions ---
% El enunciado ajusta la nota de cada miembro (±20%) según el historial de
% git y la evaluación entre pares. Esta tabla debe coincidir con los commits.
\section*{Team Contributions}

\begin{table}[H]
\centering
\caption{Contribution of each member.}
\label{tab:contributions}
\begin{tabular}{p{0.28\columnwidth}p{0.62\columnwidth}}
\toprule
Member & Main contributions \\
\midrule
Daniela Méndez Ramírez & \todo{modules, tests, report sections} \\
Nguyen Huy An Khanh & \todo{modules, tests, report sections} \\
\todo{Member 3} & \todo{modules, tests, report sections} \\
\bottomrule
\end{tabular}
\end{table}

% --- AI Tool Disclosure ---
% El enunciado exige que el material externo se cite y que cada línea sea
% entendida por quien la sube. Esta sección debe describir todos los usos.
\section*{AI Tool Disclosure}
The AI assistant Claude (Anthropic)~\cite{claude} was used in this project for the following purposes:
\begin{itemize}
    \item to improve the English writing of this report;
    \item to explain operating-system concepts, such as the scheduling metrics, the Clock policy and the deadlock detection algorithm;
    \item to help the team read and understand the scaffold, a program split into many source and header files, and how its modules connect through function pointers;
    \item to support the implementation and testing of parts of the simulator (CPU schedulers, replacement and disk policies, the workload generator and the experiment scripts), and to compare the results with our hand-computed traces.
\end{itemize}
The team reviewed and tested all the code and all the results, and takes full responsibility for them.


% --- References ---
\begin{thebibliography}{00}
\bibitem{silberschatz} A. Silberschatz, P. B. Galvin, and G. Gagne, \textit{Operating System Concepts}, 10th ed. Hoboken, NJ, USA: Wiley, 2018.
\bibitem{brief} College of Engineering and Computer Science, VinUniversity, ``COMP2040 Operating Systems -- Group Project and Suggested Themes,'' Fall 2026.
\bibitem{repo} N. H. A. Khanh, D. Méndez Ramírez, and \todo{Member 3}, ``project-1-operating-system: the \texttt{ossim} kernel simulator,'' GitHub repository, 2026. [Online]. Available: \url{https://github.com/Ankhanh975/project-1-operating-system}
\bibitem{claude} Anthropic, ``Claude,'' large language model. Accessed: Nov. 2026. [Online]. Available: \url{https://claude.ai}
\end{thebibliography}

% ==========================================
% APPENDIX
% ==========================================
\clearpage
\appendix
% --- Appendix ---
\section{Build and Run}
\label{app:commands}
The simulator and the generator are built and tested on Linux, macOS or WSL (the Linux environment of Windows):

\begin{lstlisting}
$ cd ossim
$ make
$ make test
45 passed, 0 skipped (not implemented yet), 0 failed
\end{lstlisting}

The experiments are run from \texttt{ossim/}. Each script writes its CSV files to \texttt{report/data/}, and the plots of Section~\ref{sec:experiments} are drawn from these files when the report is compiled.

\begin{lstlisting}
$ sh experiments/exp1_io.sh       # Experiment 1
$ sh experiments/exp2_frames.sh   # Experiment 2
\end{lstlisting}

One run of each experiment, for one seed:
\begin{lstlisting}
# Experiment 1 (P = 50)
$ ./wlgen --procs 10 --gap 3 --bursts 4 --cpu 6 --io-len 6 --disk 50 --io 50 --seed 1 > w.txt
$ ./ossim w.txt --sched srtf --quantum 4 --no-paging

# Experiment 2 (32 frames)
$ ./wlgen --procs 8 --pages 16 --locality 70 --io 10 --cpu 30 --bursts 4 --gap 2 --seed 1 > w.txt
$ ./ossim w.txt --sched rr --quantum 4 --disk fcfs --max-ticks 20000 --frames 32 --replace lru
\end{lstlisting}

\section{Hand-Computed Gantt Charts}
\label{app:gantt}
These are the Gantt charts of \texttt{arrivals.txt}, run with \texttt{--no-paging}. Each interval \texttt{a-b} means that the process used the CPU from tick $a$ until just before tick $b$. P1 to P6 are the jobs J1 to J6.

\begin{lstlisting}
SJF:
P1 0-7 | P4 7-9 | P3 9-12 | P1 12-16 | P5 16-21 |
P3 21-24 | P6 24-30 | P2 30-39

SRTF:
P1 0-7 | P4 7-9 | P3 9-12 | P1 12-16 | P5 16-17 |
P3 17-20 | P5 20-24 | P6 24-30 | P2 30-39

PRIO (q=2):
P1 0-2 | P2 2-11 | P1 11-12 | P6 12-18 | P4 18-20 |
P1 20-24 | P3 24-26 | P5 26-27 | P1 27-31 |
P3 31-32 | P5 32-36 | idle 36-37 | P3 37-40

MLFQ (q=2, boost after tick 19):
P1 0-2 | P2 2-4 | P3 4-6 | P4 6-8 | P5 8-10 |
P1 10-12 | P6 12-14 | P2 14-18 | P3 18-19 |
P5 19-22 | P1 22-24 | P6 24-26 | P2 26-28 |
P3 28-31 | P1 31-32 | P6 32-34 | P2 34-35 | P1 35-39
\end{lstlisting}

\section{Hand Trace of the Clock Policy}
\label{app:clock}
Table~\ref{tab:clock} traces Clock with 3 frames (F0, F1, F2) on the textbook reference string. A star (*) means that the reference bit of the frame is set, and $h$ is the frame where the hand points after the reference. When a fault happens and there is no free frame, the hand clears the bits it finds set and stops at the first frame whose bit is clear.

\begin{table}[H]
\centering
\caption{Clock with 3 frames: 14 page faults.}
\label{tab:clock}
\begin{tabular}{rrcccc c}
\toprule
\# & Ref & F0 & F1 & F2 & Fault & $h$ \\
\midrule
1 & 7 & 7* & -- & -- & yes & 0 \\
2 & 0 & 7* & 0* & -- & yes & 0 \\
3 & 1 & 7* & 0* & 1* & yes & 0 \\
4 & 2 & 2* & 0 & 1 & yes & 1 \\
5 & 0 & 2* & 0* & 1 & & 1 \\
6 & 3 & 2* & 0 & 3* & yes & 0 \\
7 & 0 & 2* & 0* & 3* & & 0 \\
8 & 4 & 4* & 0 & 3 & yes & 1 \\
9 & 2 & 4* & 2* & 3 & yes & 2 \\
10 & 3 & 4* & 2* & 3* & & 2 \\
11 & 0 & 4 & 2 & 0* & yes & 0 \\
12 & 3 & 3* & 2 & 0* & yes & 1 \\
13 & 0 & 3* & 2 & 0* & & 1 \\
14 & 3 & 3* & 2 & 0* & & 1 \\
15 & 2 & 3* & 2* & 0* & & 1 \\
16 & 1 & 3 & 1* & 0 & yes & 2 \\
17 & 2 & 3 & 1* & 2* & yes & 0 \\
18 & 0 & 0* & 1* & 2* & yes & 1 \\
19 & 1 & 0* & 1* & 2* & & 1 \\
20 & 7 & 0 & 7* & 2 & yes & 2 \\
21 & 0 & 0* & 7* & 2 & & 2 \\
22 & 1 & 0* & 7* & 1* & yes & 0 \\
\bottomrule
\end{tabular}
\end{table}

For example, at reference 4 (page 2) all three bits are set. The hand starts at F0 and clears the bits of F0, F1 and F2. It comes back to F0, whose bit is now clear, so page 7 in F0 is the victim, and page 2 is loaded there.


\end{document}
